\subsection{Trace bounds independent of dimension}\label{sec:B5T04}
For $t(X)=\sum_{ab}|X_{ab}|$, we have
\[
 t(AXB)\le\|A\|_1t(X)\|B\|_\infty.
\]
Let $D=S_*+B_*$ and $H=I+K_*$ be symmetric, with
\[
 \|S_*\|_\infty=O(k^{-1}),\ t(S_*)=O(1),\quad
 \|B_*\|_{\max}=O(k^{-2}),\ t(B_*)=O(1),
\]
\[
 \|B_*\|_\infty=O(k^{-1}),\quad
 \|K_*\|_{\max}=O(k^{-1}),\quad\|H\|_\infty=O(1).
\]
Then $t(S_*^2)=O(k^{-1})$ and
\[
 |\operatorname{tr}K_*S_*^2|=O(k^{-2}),\quad
 \|K_*S_*\|_{\max}=O(k^{-2}),\quad t(K_*S_*)=O(1).
\]
Hence $\operatorname{tr}(HS_*HS_*)=\operatorname{tr}S_*^2+O(k^{-2})$.
The terms involving $B_*$ are bounded by its largest absolute entry times
$t(HS_*H)$ or $t(HB_*H)$. We obtain
\[
 \boxed{\operatorname{tr}(HDHD)=\operatorname{tr}S_*^2+O(k^{-2}).}
\]
Moreover, $\|HD\|_{op}=O(k^{-1})$ and $\|HD\|_{HS}^2=O(k^{-1})$.
The inequality $|\operatorname{tr}T^r|\le\|T\|_{op}^{r-2}\|T\|_{HS}^2$
gives a cubic remainder of order $O(k^{-2})$ in $\log\det(I+HD)$.

\paragraph{Application to the change of support.}
In the defect $D_*$ with the exact gauge of Section~\ref{sec:B4T08}, $S_*$ contains
the entries $-p_i$ from the squares $i\in F\setminus\{c\}$.
Each of these squares meets one removed line and three retained lines, so it contributes six
oriented entries. The uniqueness of intersections prevents overlaps.
Thus
\[
 \operatorname{tr}S_*^2=6\sum_Fp_i^2+O(k^{-2}).
\]
The remaining terms come from $y$, the rank-four Schur term, and the change in the
RC projector. They satisfy the bounds for $B_*$ by Sections~\ref{sec:B5T02} and~\ref{sec:B5T03}.
The inverse $H_*$ equals $(\Gamma+\Pi)_{II}$ and satisfies the bounds for $H$.
Thus even the coefficient six is preserved, beyond the order of the term.

\paragraph{Application to the prefactor.}
Let $E=S-S_0=A^T\operatorname{diag}(3p^2+12p^3+O(p^4))A$.
We have $t(E)=O(1)$, $\|E\|_\infty=O(k^{-1})$,
$E_{ab}=O(k^{-2})$ for $a\ne b$, and
\[
 E_{aa}=3m_{2,a}+r_a,\quad\sum_a|r_a|=O(k^{-1}),
 \quad\max_a|r_a|=O(k^{-2}).
\]
Separating the diagonal and off-diagonal parts gives
\[
 \operatorname{tr}(\Gamma E\Gamma E)
 =9\sum_a m_{2,a}^2+O(k^{-2}),\qquad
 \text{cubic remainder}=O(k^{-2}).
\]
For the parent--child change, extend $m'_2$ by zero on the removed lines.
Counting each square at most four times gives
\[
 \sum_a|\Delta m_{2,a}|\le
 4\sum_Fp_i^2+4\sum_{E'}(2p_i|y_i|+y_i^2)+\sum_{a\in J}m_{2,a}
 =O(k^{-1}).
\]
Since $\max m_2,\max m'_2=O(k^{-1})$, we obtain
$\Delta\sum_a m_{2,a}^2=O(k^{-2})$.

\paragraph{Entrywise resolvent.}
With $\Delta_0=S'_{0}-S_{0,II}$, whose diagonal is zero and whose infinity norm is $O(k^{-1})$,
the Schur complements and exact gauges give
\[
 \Gamma'-\Gamma_{II}=-\Delta_0+O_{\max}(k^{-2}).
\]
Replacing the inverses by $I$ in the linear term incurs an error of
$\|H-I\|_{\max}\|\Delta_0\|_1\|H'\|_1=O(k^{-2})$.
The row norms give the same order for the quadratic terms.
Between lines of a surviving square $i$ we have $\Delta_{0,ab}=y_i$.
No removed square contributes. Therefore
\[
 \kappa'_i-\kappa_i=O(k^{-2}),\quad
 \kappa_i=a_i^T\Gamma a_i-4,
 \quad\Delta\sum_i p_i^2\kappa_i=O(k^{-2}).
\]
We also have $\Delta\sum p_i^3=O(k^{-2})$ by $\sum_Fp=O(1)$ and $y/p=O(k^{-1})$.
These are the substitutions of Sections~\ref{sec:B4T07} and~\ref{sec:B4T08}. At no point is a uniform error
multiplied by the number of diagonals.

\paragraph{Scope of the size bounds.}
The Newton remainder of relative order $k^{-2}$ has mass $O(k^{-1})$.
The linear traces of order $k^{-1}$ are retained in the cancellation of Section~\ref{sec:B5T05}.
These size bounds do not allow them to be discarded.
